% ============================================================================
% Section 3 -- the capacity model. Contribution 2.
%
% COMPRESSED 2026-09-27. A 4-page analytical section in a 10-page systems paper
% pushed the evaluation to p11; the norm puts it near p7. The derivations,
% fits, residual structure and sensitivity now live in appendix A. What stays
% is what a reader needs to USE the model: the inputs it takes, the two
% conditions it applies, and what it returns.
% ============================================================================

\section{When does a second tier pay?}
\label{sec:model}

The model takes a profile of the device and returns the host cohort that maximises
decode throughput --- or zero, when this configuration should not use a second tier
at all. Below are the profile it needs, the two conditions it applies, and what its
structure says about which configurations are worth serving this way.

\subsection{Setup}

Let $\Xlink$ be the achieved interconnect bandwidth, $\HBM$ the device memory, $W$
the weight bytes, and $\Ckv = \HBM - W - \mathcal{R}$ the memory left for KV after a
$4$\,GB device reserve for activations and fragmentation. Let $S$ be the context
length and $\btok$ the KV bytes per token.

A decode step reads every weight once and every resident KV byte once, and the two
reads do not achieve the same bandwidth. Writing $\BW_w$ and $\BW_{kv}$ for each,
%
\begin{equation}
\label{eq:step}
\mathrm{ITL}(B) \;=\; \frac{W}{\BW_w} \;+\; \frac{B \cdot S \cdot \btok}{\BW_{kv}}.
\end{equation}
%
That weights stream faster than paged KV is unsurprising; we separate them
because the gap is large enough to change the answer (\S\ref{sec:model:params}).

The capacity-limited batch is $B_{\mathrm{cap}} = \Ckv/(S\btok)$. A step of length
$\mathrm{ITL}$ cannot carry more than $\mathrm{ITL}\cdot\Xlink$ bytes across the
link, so the host cohort is bounded by
%
\begin{equation}
\label{eq:boff}
B_{\mathrm{off}} \;=\; \frac{\mathrm{ITL} \cdot \Xlink}{S \btok}.
\end{equation}
%
That is a capacity bound and nothing more: it says how many riders' bytes fit, not
what they cost. What they cost is two numbers we measure --- a \textbf{fixed} $F$
for using the host path at all, whatever the rider count, and a
\textbf{per-request} $c$.

$c$ is \emph{not} a share of the rider's transfer, and getting that wrong cost us
two rounds of this paper. Below the knee (\S\ref{sec:model:knee}) a rider's bytes
hide \emph{completely}: removing them while keeping every launch and every host
bookkeeping step moves $c$ by $0.010$\,ms, and what remains does not vary with
context --- $0.532$\,ms at 128 tokens against $0.488$ at 256. So $c$ is exposed
work only, it scales with depth rather than with bytes or bandwidth, and
\begin{equation}
\label{eq:coeffs}
c(L) \;=\; 0.177 + 0.0054\,L \ \mathrm{ms},
\qquad
F(L) \;=\; 0.053\,L \ \mathrm{ms},
\end{equation}
measured on A100/PCIe strictly below the knee: $c = 0.61$\,ms and $F = 4.24$\,ms at
$L = 80$ (\S\ref{sec:eval:cost}). $F$ carries $\pm 0.5$\,ms, which is its
run-to-run drift rather than a fit error.

\subsection{Two consequences that run against intuition}
\label{sec:model:cancel}

\paragraph{Context length cancels.} The ratio of riders to residents,
%
\begin{equation}
\label{eq:cancel}
\frac{B_{\mathrm{off}}}{B_{\mathrm{cap}}}
  \;=\; \frac{\mathrm{ITL}\cdot \Xlink}{\Ckv},
\end{equation}
%
is independent of $S$. Long context does not make offloading more attractive: what
it buys a rider in freed HBM it spends again in per-step transfer. This is worth
stating because the intuition runs the other way --- long context is what makes KV
expensive, so it looks like the regime where relocating it should pay most.

Equation~\ref{eq:step} claims ITL depends on batch and context only through their
product, and that the intercept is the weight term. On an A100-80 over 41
configurations both hold \emph{within} each model ($R^2 = 0.975$ and $0.92$), and
the intercept --- which the equation predicts rather than fits --- comes out exactly
for Qwen2.5-32B and 20\% high for Llama-3.1-8B; Appendix~\ref{app:params} gives the
fits and why the 8B misses. Pooling the models is meaningless because the intercept
is model-dependent, which is next.
% src: paper/model/p4_llama8b_a100.json, p4_qwen32b_a100.json @ 2bfc6ca5e

\paragraph{The step time is model-dependent.}
\label{sec:model:modeldep}
At capacity a step reads all of device memory once, which invites the
simplification $\mathrm{ITL}_0 \approx \HBM/\BW$ --- a quantity fixed by the device
alone. It does not hold. Evaluating Eq.~\ref{eq:step} at each model's own pool
gives $78.4$\,ms for Llama-3.1-8B against $53.2$\,ms for Qwen2.5-32B, a
$1.47\times$ spread on identical hardware, because a model that is mostly weights
reads device memory faster than one that is mostly KV --- so the gate below is
model-dependent rather than a property of the device. Both are
\emph{projections}, the largest pools fitting on an 80\,GB part measuring $69.7$
and $50.9$\,ms.
% src: paper/model/p4_llama8b_a100.json @ 2bfc6ca5e
% src: paper/model/p4_qwen32b_a100.json @ 2bfc6ca5e

\subsection{The gain law}
\label{sec:model:gain}

Adding $B_{\mathrm{off}}$ riders to $B_{\mathrm{cap}}$ residents multiplies the
served batch by
%
\begin{equation}
\label{eq:gain-exact}
\text{gain} \;=\; 1 + \frac{B_{\mathrm{off}}}{B_{\mathrm{cap}}}
      \;=\; 1 + \frac{\mathrm{ITL}_0 \cdot \Xlink}{\Ckv},
\end{equation}
%
exact in the fluid limit, and the form to compute with. Its structure is clearer on
noting that a full pool takes about $\HBM$ over the bandwidth the device achieves
across it, so the gain is roughly $(\Xlink/\BW_{\mathrm{pool}})\cdot(\HBM/\Ckv)$.
The bandwidth ratio is a handicap --- $419$\,GB/s against a pool reading at
$2000$--$2900$ on a GH200 --- that nothing a serving system improves. The amplifier
is $\HBM/\Ckv$: \textbf{the value of a second tier tracks how much of memory the
weights ate}, not model size and not context length --- the section's most useful
consequence, because it runs counter to how the field picks evaluation models.

Equation~\ref{eq:gain-exact} counts requests the device can \emph{hold}, not
requests per second it serves. Each rider also displaces a share of a resident's
blocks for the window it streams into, which is why Table~\ref{tab:gain}'s
capacity column ($\times1.50$ for Qwen2.5-32B fp16) sits below a hand
recomputation ($\times1.54$).

\paragraph{Weight precision is the clean lever.} Comparing models of different
sizes moves the weight footprint but also layer count, KV geometry and
architecture; the clean instrument is one model at two precisions, and the two
Qwen2.5-32B rows differ in nothing else. \textbf{Neither factor of the multiplier
is measured}: the bandwidth ratio needs two interconnect generations and we have
one, and the amplifier needs that precision pair run end to end. The rows are
projections.

\begin{table*}[t]
\centering
\small
\begin{tabular}{lrrrrrrr}
\toprule
Model (GH200-96) & weights & $\Ckv$ & $\HBM/\Ckv$ & ITL$_0$ & $F$ & capacity & tput \\
\midrule
Llama-8B fp16          & 16.1\,GB & 75.9\,GB & 1.26 & 46.6\,ms & \phantom{0}1.7\,ms & $\times$1.22 & $\times$1.04 \\
Qwen-32B \textbf{fp8}  & 32.0\,GB & 60.0\,GB & 1.60 & 43.3\,ms & \phantom{0}3.4\,ms & $\times$1.28 & $\times$1.08 \\
Qwen-32B \textbf{fp16} & 65.5\,GB & 26.5\,GB & 3.63 & 33.9\,ms & \phantom{0}3.4\,ms & $\times$1.50 & $\times$\textbf{1.25} \\
Llama-70B fp8          & 72.7\,GB & 19.3\,GB & 4.97 & 33.0\,ms & \phantom{0}4.2\,ms & $\times$1.67 & $\times$\textbf{1.37} \\
\bottomrule
\end{tabular}
\caption{\emph{Projections}, not measurements, at 2048 context and an 80\,ms target
--- the first round figure above $\mathrm{ITL}_0 + F$ for every row.
$\Ckv = \HBM - W - 4$\,GB. \emph{capacity} charges each rider the blocks it
displaces and the two HBM passes its KV makes; \emph{tput} additionally charges $F$
and the per-request cost. \textbf{Every row now pays, and the margin grows with
$\HBM/\Ckv$} --- the amplifier still doing the work, but on the size of the win
rather than on its sign. The slack against the break-even fixed cost $F^{*}$ runs
$2.3$, $4.3$, $10.1$ and $14.7$\,ms down the table.
\textbf{The two Qwen-32B rows differ only in weight precision.} Weights are
measured checkpoints except Qwen-32B fp8; Appendix~\ref{app:params} gives the
inputs and the sensitivity.
\textbf{The 8B row changed sign} when $F$ was re-measured --- it was
$\times 0.98$ on a fixed cost of $6.2$\,ms --- and it is the one row whose
\emph{measured} verdict is unchanged, because on A100/PCIe it admits no riders at
all (\S\ref{sec:eval:negative}).}
% src: paper/model/tables.py (wraps killer_v2.py:analyze, S=8192, GH200-96 at
% the measured 419 GB/s, slo=80 ms, stage=2, F and c from
% paper/model/overhead_model.py FINAL, measured below the knee)
\label{tab:gain}
\end{table*}

\subsection{From capacity to throughput: the second gate}
\label{sec:model:omits}

A serving system is paid in requests \emph{per second}. Charging the host path's
two costs gives
%
\begin{equation}
\label{eq:gain-tput}
\text{gain}_{\text{tput}} \;=\;
  \Big(1 + \tfrac{B_{\mathrm{off}}}{B_{\mathrm{cap}}}\Big)
  \cdot \frac{\mathrm{ITL}_0}{\mathrm{ITL}_0 + F + \Delta},
\end{equation}
%
where $\Delta = c\,B_{\mathrm{off}}$. This is a \textbf{second gate} and it binds
harder than the latency one: the technique pays only when
$B_{\mathrm{off}}/B_{\mathrm{cap}} > (F + \Delta)/\mathrm{ITL}_0$.

\textbf{All four of Table~\ref{tab:gain}'s rows clear it.} Charging each the $F$ it
would pay ($0.053$\,ms per layer, \S\ref{sec:impl:cost}), capacity gains of
$\times$1.22 to $\times$1.67 become throughput ratios of $\times$1.04,
$\times$1.08, $\times$1.25 and $\times$1.37 --- the ordering the amplifier
predicts, now on the size of the win rather than on its sign. The useful way to
report the headroom is to invert Eq.~\ref{eq:gain-tput} for the fixed cost each
configuration can afford:
%
\begin{equation}
\label{eq:fstar}
F^{*} \;=\; \Big(\tfrac{B_{\mathrm{tot}}}{B_0} - 1\Big)\mathrm{ITL}_0 - \Delta.
\end{equation}
%
Eq.~\ref{eq:fstar} gives $F^{*} = 13.4$\,ms for Qwen2.5-32B fp16 against the
$3.4$ it pays, and $19.0$ against $4.2$ for Llama-70B fp8 --- slack of $10.1$ and
$14.7$\,ms. Every row clears its threshold, and the slack widens monotonically
with $\HBM/\Ckv$: $2.3$\,ms for Llama-3.1-8B, $4.3$ for Qwen2.5-32B fp8. That is a
different claim from the one this paper made when $F$ was $12.3$\,ms, where the
same two leading rows cleared by $0.5$ and $2.8$\,ms and the other two missed.
\textbf{What changed is not the model but $F$}, and \S\ref{sec:eval:cost} reports
why the earlier figure was too large by a factor of three.

Four smaller terms are omitted --- host-side writes, host DRAM bandwidth, the
integrality of both cohorts, and MoE weight reads --- and \S\ref{sec:discussion}
states which regimes each of them bounds.

\subsection{The knee, and why the gain is a sawtooth in context}
\label{sec:model:knee}

\emph{Question: how many riders may a step admit before the transfer stops hiding?}
Eq.~\ref{eq:gain-tput} is only valid while the step is still compute-bound. Once a
rider's bytes take longer to move than the step takes to compute, the step
\emph{is} the transfer and no amount of overlap helps. The crossover is
%
\begin{equation}
\label{eq:knee}
r^{*} \;=\; \frac{\mathrm{ITL}_0 \cdot \Xlink}{S \btok},
\end{equation}
%
the same ratio as Eq.~\ref{eq:boff} but read as a rate rather than a capacity. We
evaluate $\mathrm{ITL}_0$ here at its weight-read-bound value, which omits the
resident cohort's own KV read and is therefore a \emph{lower} bound on the compute
window a transfer can hide inside; the rider counts below are correspondingly
conservative, and Table~\ref{tab:gain}, which charges the full step, admits more.
Two things follow, and together they decide where the feature should be run.

\textbf{A faster link buys the knee, not a cheaper rider.} $c$ contains no
bandwidth term at all (Eq.~\ref{eq:coeffs}), so an interconnect cannot make a
rider cost less. What it does is move $r^{*}$ in proportion to $\Xlink$, admitting
more riders at the same per-rider price. This is the opposite of how we first
reported it, and it is why a bandwidth sweep that holds $r$ fixed finds the link
nearly irrelevant.

\textbf{Riders are integers, so the achievable gain is a sawtooth in $S$.} A step
admits $\lfloor r^{*} \rfloor$ riders, so gain peaks just below each integer
crossing of Eq.~\ref{eq:knee} and drops by a whole rider's worth just above one.
The effect is not small: for Llama-70B fp8 on GH200, 16\,384 tokens of context give
$r^{*} = 1.89$ and therefore \emph{one} rider at $\times 1.07$, while 15\,493
tokens give $r^{*} = 2.00$, two riders, and $\times 1.25$ --- an $18\%$ gain
improvement for a $5\%$ reduction in context. A model that evaluates
Eq.~\ref{eq:gain-tput} at a fractional $r^{*}$ smooths this away and will choose
the wrong context. (Qwen2.5-72B fp8, which this table does not carry, does better
still: $\times 1.12$ against $\times 1.37$, a $22\%$ improvement for $1.2\%$ less
context.)

The practical caveat is larger than the effect. $r^{*}$ depends on
$\mathrm{ITL}_0$, which we compute rather than measure on the target part, and a
$2\%$ error in it moves the crossing enough to lose the rider --- $\times 1.25$
collapses to $\times 1.05$. So the recommendation is not to sit just under a
crossing but to leave margin: targeting $r^{*} \approx 2.2$ gives $\times 1.21$ and
holds both riders through a $9\%$ error in $\mathrm{ITL}_0$.

\paragraph{$\mathrm{ITL}_0$ checks out against a published GH200 measurement.}
Our step-time model is $\mathrm{ITL}_0 = (W + \text{KV read})/(\mathrm{BW} \cdot
\eta_w)$ with $\eta_w = 0.75$ taken from A100 weight reads, and we have no
Grace-Hopper part to test it on. BOOST reports one~\cite{boost}: Llama-3.3 70B at
$1$K/$8$K context on a GH200, HBM-only, batch $6$ in $81.6$\,GB, measured
\textbf{TPOT $= 27.1$\,ms}. Reconstructing that configuration --- FP8 weights plus
an FP16 KV cache over a mean sequence of $5$K tokens, which recovers their stated
capacity to $1.5\%$ --- our expression gives $\mathbf{26.8}$\,\textbf{ms}, a
$1.1\%$ error, and implies an achieved HBM efficiency of $0.742$ against the
$0.75$ we assume. The KV precision is not stated in their table and we infer it
from the capacity arithmetic, so this is a cross-check rather than a validation;
but it is the only published GH200 decode step time for a 70B-class model we could
find, and our model lands on it.

It also settles which step time the knee should use. The same comparison puts the
\emph{weights-only} figure at $23.5$\,ms, $13\%$ below the measurement --- so
omitting the resident KV read, as we do above, understates the hideable window, and
the rider counts in this subsection are conservative by about that much.

The residual error bar then cuts one way only, and in our favour. If real
efficiency is \emph{below} $0.75$, $\mathrm{ITL}_0$ rises, which both raises $r^{*}$
and shrinks $F/\mathrm{ITL}_0$. At the lowest GH200 HBM throughput anyone reports
--- $1.9$ of $4$\,TB/s, measured by DirectKV for a zero-copy attention
kernel~\cite{directkv} --- the gain would rise to $\times 1.43$ on five riders. We
keep $0.75$ because it is both the measured value and the conservative end.
% src: paper/model/operating_point.py (selftest pins each figure in this
% paragraph; coefficients from overhead_model.py FINAL)

\subsection{The applicability gate}
\label{sec:model:gate}

Capacity binds only when the latency target exceeds the full-pool step time
\emph{plus} the fixed cost of using the host path:
%
\begin{equation}
\label{eq:gate}
\mathrm{SLO} \;>\; \frac{W}{\BW_w} + \frac{\Ckv}{\BW_{kv}} + F.
\end{equation}
%
Below it the scheduler is already refusing requests HBM could hold, because
admitting them would break latency: the system is latency-limited rather than
capacity-limited, and offloading is strictly harmful. Its effect is continuous
rather than binary --- a target only slightly above $\mathrm{ITL}_0 + F$ leaves room
for very few riders whatever the link delivers, which is why Table~\ref{tab:gain} is
computed above that sum for every row. It scopes the technique to
\emph{throughput-oriented} serving and excludes interactive chat, and also excludes
latency-sensitive agentic pipelines, which are often mis-filed as throughput work
because their requests are long even when a user is waiting.

\paragraph{Measured parameters.}
\label{sec:model:params}
Both bandwidths are measured on A100-80 across 41 configurations on two models;
$\BW_w$ is stable within a model at about $0.75$ of the measured $1935$\,GB/s peak.
$\BW_{kv}$ is \textbf{not} constant: it degrades by $1.6\times$ as the KV working
set grows, and worst exactly at the capacity-bound operating point this paper is
about --- so \textbf{a roofline charging KV at one fixed bandwidth is optimistic
precisely where it matters}. We measure that curve on Llama-3.1-8B alone, and
Table~\ref{tab:gain} applies it to working sets of 19--76\,GB; the curve, the fits
and the sensitivity are in Appendix~\ref{app:params}.

\paragraph{A published result the model is consistent with.}
\label{sec:model:directkv}
DirectKV~\cite{directkv} reports its host-KV attention slightly slower than an
HBM-resident baseline on Llama-3.1-8B and winning on OPT-30B, whose 60\,GB of
weights triple the amplifier. We put it no more strongly than \emph{consistent
with}: Appendix~\ref{app:directkv} gives three objections, including that OPT-13B's
amplifier is close to the 8B's yet they win on it.
